\chapter{Theoretical background}
\startcontents[chapters]

\section{Grammatical gender in Polish}
There are multiple options of defining grammatical genders in Polish. The one chosen for this study was proposed by Mańczak (1956), where he differentiates between 5 genders:
\begin{itemize}
\item masculine personal
\item masculine animate
\item masculine inanimate
\item feminine
\item neuter
\end{itemize}
Saloni (1976) proposed an even more elaborate system, where he talked about 9 genders (instead of one neuter there were two and additionally there were three genders for plurale tantum nouns). Such a system is, however, not necessary for the current analysis, therefore we will adopt the 5 genders proposed by Mańczak. For ease of reference we will use for them the names that Saloni used, which are: M1 (masculine personal), M2 (masculine animate), M3 (masculine inanimate), F (feminine) and N (neuter).

The difference between the genders can be observed when looking at appropriate forms of adjectives in the accusative case. Adjectives agree their genders with nouns and the accusative case provides the most diversity in terms of gender forms.

The example presented here comes from Mańczak (1956). There will be 5 nouns and one adjective for all of them used as an example: \textsl{mężczyzna} (meaning 'man'), \textsl{pies} (meaning 'dog'), \textsl{stół} (meaning 'table'), \textsl{kobieta} (meaning 'woman'), \textsl{dziecko} (meaning 'child') and \textsl{dobry} (meaning 'good').

\begin{tblr}
  {
    hlines, vlines,
    cell{1}{1} = {r=2}{l},
    cell{2}{3} = {r=4}{l},
  }
  dobrego & mężczyznę & dobrych & mężczyzn & masculine personal \\
  dobrego & psa & dobre & psy & masculine animate \\
  dobry & stół & dobre & stoły & masculine inanimate \\
  dobrą & kobietę & dobre & kobiety & feminine \\
  dobre & dziecko & dobre & dzieci & neuter \\
\end{tblr}

\section{Previous work on hybrid agreement}

\section{Polish nouns with hybrid gender}
\subsection{Socially gendered nouns}
\subsection{Profession nouns}
\subsection{Depreciative forms of M1 nouns}

\section{Hypothesis}
